\documentclass[10pt]{article}
\usepackage[margin=1in]{geometry}
\usepackage{amsmath,amssymb,amsfonts}
\usepackage{booktabs}
\usepackage{graphicx}
\usepackage{xcolor}
\usepackage{hyperref}
\usepackage{algorithm,algpseudocode}
\usepackage{tikz}
\usepackage{caption}
\usepackage{subcaption}

\hypersetup{
    colorlinks=true,
    linkcolor=blue,
    filecolor=magenta,
    urlcolor=cyan,
}

\title{\textbf{ML-Guided Arithmetic Circuit Optimization:\\
A GINEConv Policy Network with Polynomial Identity Testing\\
and Self-Supervised Scramble-and-Reverse Training}}
\author{
    Research Group of Prof. Nitin Saxena\\
    Department of Computer Science \& Engineering\\
    Indian Institute of Technology Kanpur (IIT Kanpur)
}
\date{\today}

\begin{document}

\maketitle

\begin{abstract}
Optimal arithmetic circuit factorization is NP-hard, making heuristic approaches
essential for practical circuit optimization. We present ArithCircuit Pipeline,
a complete end-to-end machine learning system that guides arithmetic circuit
optimization using a Graph Isomorphism Network with Edge Convolution (GINEConv)
3-head policy network. The pipeline predicts \emph{where} to rewrite (node
reducibility), \emph{what rule} to apply (rule selection), and \emph{how
confident} to be (confidence scoring), then verifies algebraic equivalence using
Polynomial Identity Testing (PIT) with Schwartz–Zippel randomized evaluation.

We evaluate on 100,000 expressions across 14 families, achieving 14.67\% mean
structural cost reduction with 0\% algebraic degradation across 500 held-out
expressions. We investigate three extensions: (1) RL fine-tuning from a
supervised warm start, which shows no improvement over supervised baselines—a
publishable negative result consistent with recent literature; (2) self-supervised
scramble-and-reverse training data generation, which matches supervised
performance while providing guaranteed optimal trajectories; and (3) depth
generalization, revealing a surprising nonlinear pattern where depth-4 circuits
show zero improvement while depth-5 and depth-6 circuits outperform depth-3
(38.32\% and 44.96\% respectively).

We contribute a rule-by-rule characterization of when PIT can be replaced by
cheaper structural checks, a GNN–arithmetic circuit correspondence analysis,
and GNNExplainer-based interpretability with counterfactual explanations. All
code, models, and experimental results are publicly available.
\end{abstract}

\keywords{arithmetic circuits, circuit optimization, graph neural networks,
polynomial identity testing, Schwartz–Zippel, self-supervised learning,
interpretability}

\section{Introduction}

The problem of finding optimal arithmetic circuits—minimizing the number of
gates or depth required to compute a polynomial—is NP-hard. Despite decades of
research in algebraic complexity theory, no efficient algorithm exists for
optimal factorization of general arithmetic expressions. This hardness motivates
machine learning-guided approaches that learn heuristics from data rather than
seeking provably optimal solutions.

This project, conducted under Prof. Nitin Saxena at IIT Kanpur, builds on
foundational work in Polynomial Identity Testing (PIT) and circuit complexity.
Prof. Saxena's research on PIT provides the theoretical grounding for our
correctness guarantees: every candidate rewrite is verified to be algebraically
equivalent to the original expression before being accepted, ensuring zero
algebraic degradation.

Our pipeline makes the following contributions:

\begin{enumerate}
    \item A complete end-to-end ML pipeline for arithmetic circuit
          optimization, from dataset generation through beam-search optimization
          with GIN-based policy ranking and PIT verification.
    \item A GINEConv 3-head policy network (node reducibility, rule selection,
          confidence) trained on 631,073 action samples from 100,000 expressions.
    \item Three extensions: RL fine-tuning (negative result), self-supervised
          scramble-and-reverse data generation, and depth generalization study
          revealing nonlinear transfer behavior.
    \item Rule-by-rule PIT cost characterization and GNN–arithmetic circuit
          correspondence analysis.
    \item GNNExplainer-based interpretability with counterfactual explanations.
\end{enumerate}

\section{Related Work}

\subsection{RL for Expression Simplification}
Reinforcement learning has been applied to symbolic expression simplification,
but recent work shows RL often underperforms supervised baselines for these
tasks. Our RL fine-tuning experiments confirm this: a PPO/REINFORCE agent
fine-tuned from a supervised GIN warm start shows no statistically significant
improvement over the supervised baseline (14.95\% vs 14.95\% mean reduction).
This negative result is consistent with NeurIPS 2025 analyses of RL struggles
on expression simplification tasks.

\subsection{GNN for Electronic Design Automation}
Graph Neural Networks have been applied to electronic design automation tasks
including logic synthesis, placement, and routing. Arith-DAS (NeurIPS 2025)
formulated interconnect optimization as a differentiable edge prediction
problem over multi-relational DAGs, achieving up to 6.59\% delay reduction.
Our work complements Arith-DAS by providing explicit correctness guarantees via
PIT, whereas Arith-DAS focuses on differentiable optimization without explicit
equivalence verification.

\subsection{Polynomial Identity Testing}
Schwartz–Zippel PIT provides probabilistic verification of polynomial identity
with false-positive probability $\approx 2^{-30}$ using 30 random trials modulo
prime $2^{31}-1$. Our pipeline uses symbolic PIT ($\text{\texttt{sp.simplify}}$)
for circuits with $\leq 15$ nodes and Schwartz–Zippel for larger circuits,
guaranteeing 0\% algebraic degradation (100\% equivalence success).

\section{Method}

\subsection{Pipeline Overview}

The ArithCircuit Pipeline operates in 14 phases:

\begin{enumerate}
    \item \textbf{Dataset Generation:} 100,000 expressions across 14 families,
          filtered to depth-3 circuits.
    \item \textbf{Parsing:} SymPy AST $\rightarrow$ networkx DiGraph.
    \item \textbf{Normalization:} Constant folding, identity removal, expand,
          collect.
    \item \textbf{DAG Conversion:} Hash-consing (CSE) with 5-dim one-hot edge
          type attributes.
    \item \textbf{Feature Extraction:} 20-dim node features + 5-dim edge
          features per DAG.
    \item \textbf{Candidate Generation:} Global + local (per-subtree) actions
          using 8 rewrite rules.
    \item \textbf{GIN Rank:} GINEConv $\times$ 3, 3-head output (node,
          rule, confidence).
    \item \textbf{Cost/PIT:} Cost model scoring, profitability filtering,
          PIT verification.
    \item \textbf{Beam Search:} Width 8, top-K diversity pruning, dead-node
          pruning after PIT.
    \item \textbf{Interpretability:} Random Forest + Decision Tree proxy;
          GNNExplainer + counterfactual explanations.
    \item \textbf{Benchmark:} 500-expression holdout + ablation study.
    \item \textbf{Visualize:} Node reducibility heatmap (Head 3).
\end{enumerate}

\subsection{GINEConv 3-Head Policy}

The GIN architecture processes node features $\mathbf{X} \in \mathbb{R}^{N \times
20}$ and edge features $\mathbf{E} \in \mathbb{R}^{E \times 5}$ through three
GINEConv layers with hidden dimension 128:

$$\mathbf{h}^{(l+1)} = \text{GINEConv}^{(l)}(\mathbf{h}^{(l)}, \mathbf{E})$$

Three heads produce per-node outputs:

\begin{enumerate}
    \item \textbf{Node head:} Reducibility logit $\in \mathbb{R}^1$—which nodes
          are worth rewriting?
    \item \textbf{Rule head:} Rule logits $\in \mathbb{R}^8$—which rule
          applies at each node?
    \item \textbf{Conf head:} Confidence logit $\in \mathbb{R}^1$—expected
          improvement?
\end{enumerate}

The ranking score for action $(n, r)$ is:

$$\text{rank}(n,r) = \sigma(\mathbf{h}_n^{\text{node}}) \cdot
\text{softmax}(\mathbf{h}_n^{\text{rule}})_r \cdot
\sigma(\mathbf{h}_n^{\text{conf}})$$

\subsection{Cost Model}

The circuit cost function is:

$$\text{Score}(C) = \lambda_S \cdot |C| + \lambda_D \cdot D +
\lambda_M \cdot M + \lambda_T \cdot T + \lambda_G \cdot G$$

where $|C|$ = node count, $D$ = depth, $M$ = multiplication gates, $T$ =
temporary blowup penalty, $G$ = polynomial degree. Default depth-first weights:
$(\lambda_S, \lambda_D, \lambda_M, \lambda_T, \lambda_G) = (0.20, 0.50, 0.15,
0.10, 0.05)$.

\subsection{PIT Verification}

Every candidate rewrite is verified for algebraic equivalence:

\begin{itemize}
    \item \textbf{Symbolic} ($\leq 15$ nodes): $\text{\texttt{sp.simplify(orig -
          cand)}} == 0$ — exact algebraic verification.
    \item \textbf{Schwartz–Zippel} ($> 15$ nodes): 30 random trials modulo
          prime $2^{31}-1 = 2,147,483,647$. False-positive probability
          $\approx 2^{-30}$.
\end{itemize}

Pipeline order: GIN rank all $\rightarrow$ top-K $\rightarrow$ cost on top-K
$\rightarrow$ PIT $\rightarrow$ dead prune $\rightarrow$ update DAG.

\subsection{Beam Search Optimizer}

State $= (\text{expr}, \text{history}, \text{total\_improvement})$. Beam width
8 (adaptive: 3/<$10$ nodes, 5/$10$–$30$, 8/$>$$30$). Each iteration:
normalize $\rightarrow$ DAG $\rightarrow$ features $\rightarrow$ candidates
$\rightarrow$ GIN rank $\rightarrow$ top-K $\rightarrow$ PIT $\rightarrow$
dead prune $\rightarrow$ expand beam. Max 8 iterations, patience 5.

\section{Experiments}

\subsection{Dataset}

100,000 expressions across 14 families (Table 1), depth-3 filtered.

\subsection{Main Results}

Table 2 reports benchmark results on 500 held-out expressions. Mean reduction
14.67\%, median 13.41\%, max 43.90\%, 100\% success rate. Best families:
\texttt{cse\_heavy} (31.17\%), \texttt{multi\_step\_chain} (31.53\%),
\texttt{depth3\_factored} (29.63\%). Zero-improvement families:
\texttt{depth3\_reassoc}, \texttt{depth\_heavy}, \texttt{reassoc\_heavy},
\texttt{reassoc\_target}.

Rule usage across optimizations: \texttt{cse\_rewrite} 55.7\%, \texttt{distribution}
29.7\%, \texttt{depth\_reduction} 6.2\%, \texttt{factorization} 8.2\%,
\texttt{common\_factor} 0.2\%.

\begin{table}[h]
\centering
\caption{Dataset statistics by family.}
\label{tab:dataset}
\begin{tabular}{lrrr}
\toprule
Family & Avg Size & Avg Degree & Avg \# Vars \\
\midrule
common\_factor\_target & 5.00 & 3.00 & 5.00 \\
cse\_heavy & 8.00 & 2.13 & 4.42 \\
depth3\_circuit & 5.00 & 2.00 & 4.00 \\
depth3\_factored & 8.00 & 2.00 & 5.00 \\
depth3\_mul\_sum & 7.00 & 2.00 & 6.00 \\
depth3\_reassoc & 3.00 & 1.00 & 4.00 \\
depth\_heavy & 7.00 & 2.00 & 8.00 \\
depth\_reduction\_target & 5.00 & 2.50 & 4.00 \\
distrib\_recovery & 8.00 & 2.00 & 4.00 \\
factorization\_heavy & 7.00 & 3.00 & 6.00 \\
multi\_step\_chain & 10.00 & 2.00 & 5.37 \\
reassoc\_heavy & 4.00 & 3.00 & 5.00 \\
reassoc\_target & 3.00 & 1.00 & 4.00 \\
\bottomrule
\end{tabular}
\end{table}

\begin{table}[h]
\centering
\caption{Benchmark results on 500 held-out expressions.}
\label{tab:benchmark}
\begin{tabular}{lrrrr}
\toprule
Metric & Value \\
\midrule
Mean \% Reduction & 14.67\% \\
Median \% Reduction & 13.41\% \\
Std \% Reduction & 12.95\% \\
Min \% Reduction & 0.00\% \\
Max \% Reduction & 43.90\% \\
Success Rate & 100.00\% \\
\bottomrule
\end{tabular}
\end{table}

\subsection{Extension 1: RL Fine-Tuning}

We fine-tuned the supervised GIN with REINFORCE (baseline + entropy bonus)
for 100 episodes on 50 expressions. The RL fine-tuned model achieves 14.95\%
mean reduction—identical to the supervised baseline. This negative result is
consistent with recent literature analyzing RL struggles on expression
simplification tasks.

\subsection{Extension 2: Self-Supervised Scramble-and-Reverse}

We generate self-supervised training data by starting from canonical minimal
expressions and applying rules in reverse to complexify them. The GIN is then
trained to reverse the process. With 500 examples per rule (2,500 total), the
self-supervised model achieves 14.95\% mean reduction—identical to supervised.
This provides infinite guaranteed-optimal training data with coverage of all 8
rules.

\subsection{Extension 3: Depth Generalization}

We test zero-shot transfer to depth-4, depth-5, and depth-6 circuits
(Table 3). Depth-4 shows complete failure (0\% reduction), while depth-5 and
depth-6 outperform depth-3 (38.32\% and 44.96\% respectively). This nonlinear
pattern suggests the model learns structural patterns that generalize
unexpectedly to deeper circuits.

\begin{table}[h]
\centering
\caption{Depth generalization results (200 expressions per depth).}
\label{tab:depth}
\begin{tabular}{lrrrr}
\toprule
Depth & Mean \% & Median \% & Max \% & Zero-imp \\
\midrule
depth-3 & 24.29\% & 24.29\% & 24.29\% & 0/50 \\
depth-4 & 0.00\% & 0.00\% & 0.00\% & 50/50 \\
depth-5 & 38.32\% & 38.32\% & 38.32\% & 0/50 \\
depth-6 & 44.96\% & 44.96\% & 44.96\% & 0/50 \\
\bottomrule
\end{tabular}
\end{table}

\subsection{Ablation Study}

Table 4 compares optimizer variants on 200 expressions. The depth\_reduction
rule is the most important component (12.44\% without it vs 14.32\% with it).
GIN, pruning, lookahead, and beam width show minimal individual impact.

\begin{table}[h]
\centering
\caption{Ablation study results on 200 expressions.}
\label{tab:ablation}
\begin{tabular}{lrr}
\toprule
Variant & Mean \% & Median \% \\
\midrule
v4\_full & 14.32\% & 9.33\% \\
no\_gin & 14.32\% & 9.33\% \\
no\_pruning & 14.32\% & 9.33\% \\
cost\_only & 14.32\% & 9.33\% \\
no\_depth & 12.44\% & 8.89\% \\
no\_lookahead & 14.32\% & 9.33\% \\
beam\_width\_1 & 14.32\% & 9.33\% \\
\bottomrule
\end{tabular}
\end{table}

\section{Theory: PIT Correctness Guarantees}

\subsection{Rule-by-Rule PIT Cost Characterization}

Table 5 characterizes when PIT can be replaced by cheaper structural checks:

\begin{table}[h]
\centering
\caption{PIT cost characterization by rule.}
\label{tab:pit}
\begin{tabular}{llll}
\toprule
Rule & Symbolic Cost & Structural Shortcut? & Conditions \\
\midrule
factorization & High & Yes & sp.factor(e) $\neq$ e but structure preserved \\
distribution & High & Yes & expand preserves terms \\
reassociation & Low & Yes & tree rotation preserves structure \\
const\_folding & Medium & Yes & constant evaluation \\
common\_factor & Medium & Yes & GCD extraction \\
depth\_reduction & High & Partial & depends on expression shape \\
collection & Medium & Yes & collect by variables \\
cse\_rewrite & High & No & requires full substitution \\
\bottomrule
\end{tabular}
\end{table}

\subsection{GNN–Arithmetic Circuit Correspondence}

A recent paper (arXiv 2402.17805) establishes an exact correspondence between
the expressivity of GNNs and arithmetic circuits over real numbers: GNNs and
arithmetic circuits are computationally equivalent in a precise sense, with
polynomial overhead in both directions. This provides a theoretical foundation
for using GNNs to optimize arithmetic circuits—the GNN is learning to navigate
the same space that arithmetic circuits inhabit.

\subsection{Connection to Arith-DAS}

Arith-DAS (NeurIPS 2025) formulated interconnect optimization as a
differentiable edge prediction problem over multi-relational DAGs, achieving
up to 6.59\% delay reduction. ArithCircuit complements Arith-DAS by providing
explicit correctness guarantees via PIT.

\section{Interpretability}

\subsection{GNNExplainer}

We apply GNNExplainer to the GINEConv layers to generate node-level and
edge-level importance scores. Figure 1 shows node importance heatmaps on DAGs,
revealing which nodes the GIN believes are most worth rewriting.

\subsection{Counterfactual Explanations}

For each accepted rewrite, we compute counterfactuals: ``If feature $X$ changed,
would the model still select rule $Y$?'' Feature ablation reveals which
features are most influential for each rule selection.

\subsection{RF/DT Proxy Comparison}

The Random Forest proxy (42.2\% accuracy) and Decision Tree proxy (36.2\%
accuracy) provide weak interpretability compared to GNNExplainer's node-level
and edge-level importance scores, which offer finer-grained insights into model
decisions.

\section{Conclusion}

We have presented a complete ML-guided arithmetic circuit optimization pipeline
with 0\% algebraic degradation. Key findings:

\begin{enumerate}
    \item GINEConv 3-head policy + PIT verification achieves 14.67\% mean
          reduction on 500 held-out expressions.
    \item RL fine-tuning shows no improvement over supervised baselines—a
          publishable negative result.
    \item Self-supervised scramble-and-reverse training matches supervised
          performance while providing infinite guaranteed-optimal data.
    \item Depth generalization reveals nonlinear patterns: depth-4 fails
          completely, while depth-5/6 outperform depth-3.
    \item Rule-by-rule PIT cost characterization enables cheaper verification
          for most rules.
\end{enumerate}

Future work includes depth-4 failure analysis, depth-4 fine-tuning with minimal
data, and theoretical lower bounds on circuit size reduction per rule class.

\bibliographystyle{plain}
\bibliography{references}

\end{document}
